% Generated draft: requires mathematical and source review.
\documentclass[11pt]{article}
\usepackage[margin=25mm]{geometry}
\usepackage{fontspec}
\setmainfont{Latin Modern Roman}
\usepackage{amsmath,amssymb,mathtools,unicode-math}
\setmathfont{Latin Modern Math}
\usepackage{enumitem,needspace,xurl,hyperref}
\setlength{\parindent}{0pt}
\setlength{\parskip}{5pt}
\setlength{\emergencystretch}{3em}
\newcommand{\N}{\mathbb{N}}
\newcommand{\Z}{\mathbb{Z}}
\newcommand{\Q}{\mathbb{Q}}
\newcommand{\R}{\mathbb{R}}
\newcommand{\C}{\mathbb{C}}
\newcommand{\F}{\mathbb{F}}
\newcommand{\A}{\mathbb{A}}
\newcommand{\PP}{\mathbb{P}}
\newcommand{\E}{\mathbb{E}}
\newcommand{\eps}{\varepsilon}
\newcommand{\dd}{\,\mathrm{d}}
\begin{document}
\section*{UnsolvedMath ID: 30005999}
\textbf{Source problem code:} catalogue entry
\subsection{Definitions and mathematical statement}
\textbf{Terminology and scope.} Individual dislocations are the microscopic defects whose physical motion laws are intended to provide the starting point. A dislocation-density field is a continuum description of a distribution of dislocations. The source does not specify its mathematical type or the information it records. The associated elastoplastic effects refer to the coupled elastic response and plastic deformation arising from dislocation motion; their state variables and constitutive relations are likewise unspecified.

\textbf{Source-stated research problem.} Derive continuum evolution laws for dislocation-density fields, together with their associated elastoplastic effects, from the physical laws governing individual dislocations. The original discussion describes the formulation of laws for dislocation-field motion and the corresponding elastoplastic effects as a pressing open problem of solid mechanics, despite substantial consensus concerning the physical laws for individual dislocations.

\textbf{Unresolved mathematical specification.} This is a research objective rather than a fully specified mathematical conjecture. The supplied record gives no microscopic equations, spatial domain or dimension, admissible dislocation configurations, interaction laws, loading, initial or boundary conditions, scaling regime, continuum state space, convergence notion, or proposed limiting evolution law. Consequently, it does not determine a proposition with fixed hypotheses and quantifiers. A precise homogenization statement would require these choices and a specified relation between microscopic trajectories and continuum evolution; none is selected here.

\textbf{Qualification concerning existing results.} The supplied dated assessment states that Bonicatto and Rindler rigorously homogenize a prototypical rate-independent model from discrete dislocations to dislocation fields, while leaving general constitutive laws for dislocation-field motion and elastoplastic effects unsettled. The assumptions and conclusions of that particular result are not supplied, so it cannot be substituted for a precise formulation of the broader target.
\subsection{Short English statement}
The problem asks how motion laws for individual dislocations lead to continuum laws for dislocation-density fields and the resulting elastic and plastic deformation. The source leaves the microscopic model, continuum variables, scaling, and initial and boundary conditions unspecified, so a unique mathematical conjecture cannot yet be stated. The supplied assessment labels the problem partially solved through a prototypical rate-independent homogenization result; that dated assessment requires independent verification.
\subsection{Sources}
\begin{itemize}
\item Local UnsolvedMath registry, permanent ID \(30005999\), title \emph{Continuum Laws for Dislocation Motion}, source problem code \texttt{catalogue entry}; supplied complete record from the Oberwolfach Reports Open Problems collection. Repository registry locations are \texttt{data/unsolvedmath/problems.json} and \texttt{catalog/problems.json}; the individual file containing this record has not been independently checked.
\item Supplied report citation: \emph{Calculus of Variations (2025)}, Oberwolfach Reports, DOI \texttt{10.4171/owr/2024/37}, pp.\ 2113--2178. \url{https://doi.org/10.4171/owr/2024/37}. The record assigns proposed year \(2024\); the supplied citation's \(2025\) designation and the report's \(2024\) identification are retained without reconciliation.
\item Attribution in the original statement: Filip Rindler, joint work with Paolo Bonicatto, \emph{On the homogenization problem for elasto-plasticity driven by dislocation motion}. The record supplies no separate named proposer.
\item Supplied literature reference: Bonicatto--Rindler, dislocation homogenization. \url{https://arxiv.org/abs/2410.02906}.
\item Supplied literature reference: Oberwolfach Report 37/2024. \url{https://ems.press/content/serial-article-files/50045}.
\item Literature assessment date: \texttt{2026-08-22}. This is supplied dated curation, not an independently verified assessment. It labels the problem \texttt{partially\_solved} and reports a rigorous homogenization result for a prototypical rate-independent model, while identifying broader constitutive laws as unresolved. The background explicitly requests independent verification of time-sensitive status and cited claims.
\item The registry labels the statement \texttt{corrected\_verified} and says it was reconstructed and checked against the official Oberwolfach report. These are supplied provenance claims; no source has been independently visited or verified for this draft.
\end{itemize}
\end{document}
